% Proposed mechanisms, with qualitative sketches of support and knot allocation.
% Four numbered stages on a shared flow line, then an unruled sketch strip.
% Each sketch carries its lettered title BELOW it, centred on the sketch's own
% horizontal extent (not on an arbitrary grid), so title and drawing read as one
% unit when the strip is scanned from left to right.
\begingroup
% Shared editable vector language for the five Section 3.2 technical routes.
%
% The figures are explanatory diagrams, not experimental results.  A single
% visual grammar is used throughout:
%   * a dark header band names a block; white type is reserved for headers;
%   * a numbered marker gives the reading order of a route, and the marker of a
%     stage that carries the study's own mechanism is filled;
%   * white cards with a thin theme-coloured rule hold stage titles and bodies;
%   * tinted panels hold mechanism sketches, and a heavier frame marks the
%     representation or relation that is handed to the next study;
%   * solid arrows are information flow, dashed arrows are observation
%     feedback or a model-predicted interval.
% Colour is always paired with a marker, line style or label so that the
% grayscale proof remains readable.
\linespread{1}\selectfont
\usetikzlibrary{arrows.meta,calc,positioning,matrix,decorations.pathreplacing}

% ---- palette -------------------------------------------------------------
\definecolor{trBlueDark}{HTML}{1F4E79}
\definecolor{trBlueMid}{HTML}{2F6690}
\definecolor{trBlueSoft}{HTML}{7FA6C4}
\definecolor{trBluePale}{HTML}{E8F1F8}
\definecolor{trBlueWash}{HTML}{F4F8FB}
\definecolor{trGreenDark}{HTML}{2E5E40}
\definecolor{trGreenMid}{HTML}{3E7A52}
\definecolor{trGreenSoft}{HTML}{8FB39A}
\definecolor{trGreenPale}{HTML}{E9F2EB}
\definecolor{trGreenWash}{HTML}{F4F9F5}
\definecolor{trOchreDark}{HTML}{7A5018}
\definecolor{trOchreMid}{HTML}{9E6B27}
\definecolor{trOchreSoft}{HTML}{C9A46B}
\definecolor{trOchrePale}{HTML}{F8EEDC}
\definecolor{trOchreWash}{HTML}{FCF7EE}
\definecolor{trVioletDark}{HTML}{4A3A6A}
\definecolor{trVioletMid}{HTML}{67528E}
\definecolor{trVioletSoft}{HTML}{A294BE}
\definecolor{trVioletPale}{HTML}{EFEAF6}
\definecolor{trVioletWash}{HTML}{F8F5FC}
\definecolor{trInkSoft}{HTML}{4D4D4D}
\definecolor{trRule}{HTML}{9AA5AE}

\newcommand{\trTheme}[1]{%
  \colorlet{trDark}{tr#1Dark}%
  \colorlet{trMid}{tr#1Mid}%
  \colorlet{trSoft}{tr#1Soft}%
  \colorlet{trPale}{tr#1Pale}%
  \colorlet{trFill}{tr#1Pale}%
  \colorlet{trWash}{tr#1Wash}%
}
\trTheme{Blue}

% ---- shared type scale for every Part 3 figure ----------------------------
% Set sizes here rather than overriding individual diagrams.  The report uses
% the natural figure width, so these remain the printed sizes; subscripts use
% the normal mathematical script sizes.
\newcommand{\trFontHead}{\songti\bfseries\fontsize{9.2pt}{11.2pt}\selectfont}
\newcommand{\trFontSubhead}{\songti\bfseries\fontsize{8.2pt}{10.2pt}\selectfont}
\newcommand{\trFontBody}{\songti\fontsize{8.2pt}{10.2pt}\selectfont}
\newcommand{\trFontSmall}{\songti\fontsize{8pt}{10pt}\selectfont}
\newcommand{\trFontTiny}{\songti\fontsize{8pt}{10pt}\selectfont}
\newcommand{\trFontMath}{\fontsize{9.2pt}{11.2pt}\selectfont}
\newcommand{\trFontMathSmall}{\fontsize{8.5pt}{10.4pt}\selectfont}
\newcommand{\trFontMarker}{\songti\bfseries\fontsize{8pt}{10pt}\selectfont}
\newcommand{\trFontSubcaption}{\songti\mdseries\fontsize{9.2pt}{11.2pt}\selectfont}

% ---- base styles ---------------------------------------------------------
\tikzset{
  trBase/.style={x=1cm,y=1cm,text=black,line cap=round,line join=round,
    font=\trFontBody,
    every node/.style={align=center,outer sep=0pt,inner sep=0pt}},
  trTitle/.style={font=\trFontHead},
  trTitleSmall/.style={font=\trFontSubhead},
  trHead/.style={font=\trFontHead,text=white},
  trBody/.style={font=\trFontBody},
  trNote/.style={font=\trFontSmall},
  trTiny/.style={font=\trFontTiny},
  trLabel/.style={font=\trFontSmall},
  trMath/.style={font=\trFontMath},
  trMathSmall/.style={font=\trFontMathSmall},
  trNum/.style={circle,draw=trDark,fill=trDark,text=white,
    line width=0.5pt,inner sep=0pt,minimum size=0.36cm,
    font=\trFontMarker},
  trNumOpen/.style={trNum,fill=white,text=trDark},
  trCard/.style={fill=white,draw=trMid,line width=0.55pt,rounded corners=1.2pt},
  trFrame/.style={draw=black,line width=0.65pt},
  trArrow/.style={draw=trDark,line width=0.8pt,
    -{Latex[length=1.7mm,width=1.15mm]}},
  trArrowThin/.style={draw=trDark,line width=0.6pt,
    -{Latex[length=1.4mm,width=0.95mm]}},
  trFlow/.style={draw=black!72,line width=0.7pt,
    -{Latex[length=1.6mm,width=1.05mm]}},
  trFeedback/.style={draw=trDark,line width=0.8pt,dash pattern=on 3pt off 2pt,
    -{Latex[length=1.7mm,width=1.15mm]}},
  trPredict/.style={draw=trMid,line width=0.75pt,dash pattern=on 2.2pt off 1.6pt,
    -{Latex[length=1.6mm,width=1.05mm]}},
  trGuide/.style={draw=black!35,line width=0.45pt},
  trTag/.style={draw=trMid,fill=white,line width=0.45pt,inner sep=2.2pt,
    rounded corners=0.8pt,font=\trFontSmall},
  trAccent/.style={draw=trMid,line width=0.9pt},
  trAccentThin/.style={draw=trMid,line width=0.6pt},
  trPossible/.style={draw=trSoft,line width=0.7pt,dash pattern=on 1.8pt off 1.4pt},
  trTimeLine/.style={draw=black!55,line width=0.6pt,
    -{Latex[length=1.5mm,width=0.95mm]}},
  trDown/.style={draw=trDark,line width=0.8pt,
    -{Latex[length=1.6mm,width=1.05mm]}},
  trDouble/.style={draw=trDark,line width=0.8pt,
    {Latex[length=1.5mm,width=1.0mm]}-{Latex[length=1.5mm,width=1.0mm]}},
  trDoubleSmall/.style={draw=trDark,line width=0.7pt,
    {Latex[length=1.1mm,width=0.75mm]}-{Latex[length=1.1mm,width=0.75mm]}}
}

% ---- reusable pieces -----------------------------------------------------
% Header band.  #1 left, #2 right, #3 bottom, #4 top, #5 text.
\newcommand{\trHeader}[5]{%
  \path[fill=trDark] (#1,#3) rectangle (#2,#4);
  \node[trHead] at ({#1/2+#2/2},{#3/2+#4/2}) {#5}
}

% ---- compact detailed-route primitives (2026-10-30) ----------------------
% The four detailed routes use a downward y axis so that a route reads from
% the input strip at the top to the output strip at the bottom.  Their canvases
% are 15.56 cm wide and only as tall as the content; the styles below keep the
% four redraws on the same type scale as the overview (Fig. 11).
\tikzset{
  trRouteBase/.style={x=1cm,y=-1cm,text=black,line cap=round,line join=round,
    font=\trFontBody,
    every node/.style={align=center,outer sep=0pt,inner sep=0pt}},
  trRouteHead/.style={font=\trFontHead,text=white},
  trRouteBody/.style={font=\trFontBody},
  trRouteSmall/.style={font=\trFontSmall},
  trRouteTiny/.style={font=\trFontTiny},
  trRouteMath/.style={font=\trFontMath},
  trRouteMathSmall/.style={font=\trFontMathSmall},
  trRouteFlow/.style={draw=trDark,line width=0.72pt,
    -{Latex[length=1.5mm,width=1.0mm]}},
  trRouteThin/.style={draw=trDark,line width=0.6pt,
    -{Latex[length=1.3mm,width=0.9mm]}},
  trRouteBack/.style={draw=trDark,line width=0.65pt,
    dash pattern=on 2.6pt off 1.8pt,-{Latex[length=1.4mm,width=0.95mm]}},
  trRouteRule/.style={draw=trMid,line width=0.55pt},
  trRouteMotion/.style={draw=trBlueDark,line width=0.95pt},
  trRoutePossible/.style={draw=black!52,line width=0.65pt,
    dash pattern=on 1.8pt off 1.4pt},
  trRouteCard/.style={fill=white,draw=trMid,line width=0.6pt},
  trRouteStrip/.style={fill=trWash,draw=trMid,line width=0.55pt},
  trRouteFocus/.style={fill=trPale,draw=trDark,line width=0.7pt}
}

\trTheme{Time}
\begin{tikzpicture}[trBase]

  \trStage{0.30}{2.95}{1}{观测线索}{二维对应、几何及可见性}
  \trStage[core]{3.95}{2.95}{2}{可靠性估计}{循环一致性、几何重投影支持}
  \trStage[core]{7.60}{2.95}{3}{时间配置}{有效作用区间、样条表达自由度}
  \trStage{11.25}{2.95}{4}{联合求解}{连续轨迹、几何及外观}

  \trStep{3.25}{3.95}
  \trStep{6.90}{7.60}
  \trStep{10.55}{11.25}

  % joint solving feeds back a geometric correction to the reliability estimate
  \draw[trFeedback] (12.72,4.30) -- (12.72,3.66) -- (5.42,3.66) -- (5.42,4.30);
  \node[trTag] at (9.05,3.66) {几何回校正};

  % ---- sketch strip: drawing above, centred title below ----
  % (a) observation support defines a local valid interval I_j   [0.60 .. 4.10]
  \path[fill=trTimeFill] (1.35,1.85) rectangle (3.35,2.82);
  \draw[trArrow] (0.65,2.10) -- (4.15,2.10);
  \foreach \xx in {0.95,1.35,1.85,2.35,2.85,3.35,3.85}{
    \draw[trSketch] (\xx,2.01) -- (\xx,2.19);
  }
  \fill (1.35,2.10) circle (0.058);
  \fill (3.35,2.10) circle (0.058);
  \draw[trMotion] (1.35,2.88) -- (1.35,3.00) -- (3.35,3.00) -- (3.35,2.88);
  \node[trLabel] at (2.35,3.12) {$I_j$};
  \node[trLabel,anchor=west] at (4.07,1.86) {$t$};
  \node[text width=3.50cm,align=center] at (2.35,1.33)
    {{\captionsetup[subfigure]{font=figsubroute,skip=0pt}
      \subcaption{局部时间支撑}\label{subfig:route-time-support}}};

  % (b) knot density follows local complexity   [5.05 .. 8.95]
  \draw[trMotion] plot[smooth] coordinates {
    (5.10,2.10) (6.00,2.24) (6.65,2.15) (7.25,2.92) (7.80,2.15) (8.90,2.30)};
  \foreach \xx/\yy in {5.10/2.10,6.00/2.24,6.65/2.15,7.01/2.60,7.25/2.92,7.50/2.62,7.80/2.15,8.90/2.30}{
    \draw[trControl] (\xx,\yy) circle (0.056);
  }
  \node[text width=3.90cm,align=center] at (7.00,1.33)
    {{\captionsetup[subfigure]{font=figsubroute,skip=0pt}
      \subcaption{表达能力配置}\label{subfig:route-time-capacity}}};

  % (c) position and velocity readable at a non-sampled time   [9.75 .. 13.10]
  \draw[trMotion] plot[smooth] coordinates {
    (9.80,2.10) (10.60,2.72) (11.40,2.48) (12.20,2.10) (13.05,2.72)};
  \fill (11.40,2.48) circle (0.058);
  \draw[trArrow,draw=trResponseAccent] (11.40,2.48) -- (12.28,2.14);
  \node[trLabel] at (11.10,2.86) {$\mathbf{x}_i(t)$};
  \node[trLabel,anchor=west] at (12.35,1.98) {$\mathbf{v}_i(t)$};
  \node[text width=3.35cm,align=center] at (11.42,1.33)
    {{\captionsetup[subfigure]{font=figsubroute,skip=0pt}
      \subcaption{连续状态查询}\label{subfig:route-time-query}}};
\end{tikzpicture}
\endgroup
